\subsection{对称张量的余积}

设 M 为自由 A-模，且 \(\Delta_{M} = \Delta\) 为对角同态 \(x \mapsto (x, x)\) （把 M 映入 \(M \oplus M\) ）。设 \(c_{M} = c\) 为由以下同态组成的分次 A-代数的保单位同态：

\[
\mathbf {T S} (\mathrm{M}) \xrightarrow {\mathbf {T S} (\Delta)} \mathbf {T S} (\mathrm{M} \oplus \mathrm{M}) \xrightarrow {\sigma} \mathbf {T S} (\mathrm{M}) \otimes \mathbf {T S} (\mathrm{M})
\]

，其中 \(\sigma\) 是典范同构。配备 \(c\) ，\(\mathbf{TS}(\mathbf{M})\) 是一个分次 A-余代数。

对所有 \(x \in \mathbf{M}\) 及每个整数 \(p \geqslant 0\) ，我们有 \(\mathbf{TS}(\Delta) (\gamma_p(x)) = \gamma_p((x, x))\) ，因此由 (8)，

\[
c \left(\gamma_ {p} (x)\right) = \sum_ {r + s = p} \gamma_ {r} (x) \otimes \gamma_ {s} (x).\tag{10}
\]

特别地

\[
c (x) = x \otimes 1 + 1 \otimes x.\tag{11}
\]

设 \((x_{i})_{i\in I}\) 是 \(\mathbf{M}\) 的一个元素族，且对于 \(\nu \in \mathbf{N}^{(I)}\) ，令 \(x_{\nu} = \prod_{i\in I}\gamma_{\nu_i}(x_i)\) 。则

\[
c \left(x _ {\nu}\right) = \sum_ {\rho + \sigma = \nu} x _ {\rho} \otimes x _ {\sigma}.\tag{12}
\]

这由(10)得出，因为 \(c\) 是一个代数同态。

命题 8. --- 设 M 是自由 A-模，则以其代数与余代数结构，TS(M) 是一个分次交换且余交换的双代数。余单位是 A-线性映射 \(\varepsilon: \mathbf{TS}(\mathbf{M}) \to \mathbf{TS}^0(\mathbf{M}) = \mathbf{A}\) ，它在 \(\mathbf{TS}^p(\mathbf{M})\) （\(p > 0\) ）上为零，且满足 \(\varepsilon(1) = 1\) 。

我们知道 A-代数 TS(M) 是结合的、交换的且含单位元的。另一方面，余积按构造是分次代数的同态；而余代数 TS(M) 是余结合且余交换的这一事实，可由公式 (10) 通过简单计算得出。映射 \(\varepsilon\) 从 TS(M) 到 A 是分次代数的同态，且满足 \(\varepsilon(1)=1\) 。最后，对每个 \(x\in M\) 我们有 \(\varepsilon(\gamma_{p}(x))=0\) ，若 p\textgreater0, \(\varepsilon(\gamma_{0}(x))=1\) ；若我们记住 (10)，这表明 \((\varepsilon\otimes1)\circ c=(1\otimes\varepsilon)\circ c=\mathrm{Id}_{\mathbf{TS}(\mathbf{M})}\) ；因此 \(\varepsilon\) 是 TS(M) 的余单位。

命题 9. --- 设 M 和 N 是自由 A-模，u 是 M 到 N 的一个 A-同态；则 TS(u) 是双代数同态。

因为我们有 \(\Delta_{\mathbf{N}} \circ u = (u, u) \circ \Delta_{\mathbf{M}}\) ，所以图表

\[
\begin{array}{c} \mathsf {T S (M)} \xrightarrow {\mathsf {T S} (\Delta_ {M})} \mathsf {T S (M \oplus M)} \xrightarrow {\sigma} \mathsf {T S (M)} \otimes \mathsf {T S (M)} \\ \mathsf {T S (u)} \Bigg | _ {\mathsf {T S (u)}} \\ \mathsf {T S (N)} \xrightarrow {\mathsf {T S} (\Delta_ {N})} \mathsf {T S (N \oplus N)} \xrightarrow {\tau} \mathsf {T S (N)} \otimes \mathsf {T S (N)}, \end{array}
\]

，其中 \(\sigma\) 和 \(\tau\) 是典范同构，该图表交换（IV，第 49 页）。因此 \(c_{\mathrm{N}} \circ \mathbf{T}\mathbf{S}(u) = (\mathbf{T}\mathbf{S}(u) \otimes \mathbf{T}\mathbf{S}(u)) \circ c_{\mathrm{M}}\) 。

命题 10. --- 设 M 是自由 A-模，则双代数 TS(M) 的本原元（III，第 602 页）恰为 M 的元素。

设 \((e_i)_{i\in I}\) 是 \(\mathbf{M}\) 的一组基，且对于 \(\nu \in \mathbf{N}^{(I)}\) ，令 \(e_{\nu} = \prod_{i\in I}\gamma_{\nu_i}(e_i)\) 。设

\(z = \sum_{\nu \in \mathbf{N}^{(I)}} \lambda_{\nu} e_{\nu}\) 为 \(\mathbf{T}\mathbf{S}(\mathbf{M})\) 的一个元素，则由 (12) 我们有

\[
c (z) = \sum_ {\nu} \lambda_ {\nu} \sum_ {\rho , \sigma \in \mathbf {N} ^ {(I)}, \rho + \sigma = \nu} e _ {\rho} \otimes e _ {\sigma} = \sum_ {\rho , \sigma} \lambda_ {\rho + \sigma} e _ {\rho} \otimes e _ {\sigma}
\]

因此

\[
c (z) - 1 \otimes z - z \otimes 1 = \sum_ {\rho \neq 0, \sigma \neq 0} \lambda_ {\rho + \sigma} e _ {\rho} \otimes e _ {\sigma} - \lambda_ {0} e _ {0} \otimes e _ {0}.
\]

因此

\[
\begin{array}{r l} z \text {   primitive   } & \Leftrightarrow \lambda_ {\rho + \sigma} = 0 \text {   when   } \rho \neq 0 \text {   and   } \sigma \neq 0 \text {   and   } \lambda_ {0} = 0 \\ & \Leftrightarrow \lambda_ {\nu} = 0 \text {   when   } | \nu | \neq 1 \\ & \Leftrightarrow z \in M. \end{array}
\]

